% Created 2023-09-12 mar. 13:17
% Intended LaTeX compiler: pdflatex
\documentclass[11pt]{article}
\usepackage[utf8]{inputenc}
\usepackage[T1]{fontenc}
\usepackage{graphicx}
\usepackage{longtable}
\usepackage{wrapfig}
\usepackage{rotating}
\usepackage[normalem]{ulem}
\usepackage{amsmath}
\usepackage{amssymb}
\usepackage{capt-of}
\usepackage{hyperref}
\date{\today}
\title{Poster Highlight speech}
\hypersetup{
 pdfauthor={},
 pdftitle={Poster Highlight speech},
 pdfkeywords={},
 pdfsubject={},
 pdfcreator={Emacs 28.2 (Org mode 9.6.1)}, 
 pdflang={English}}
\usepackage{biblatex}
\addbibresource{~/Documents/studies/phd/manuscript/references/references.bib}
\begin{document}

\maketitle
Plan
    Contexte
        -> On veut générer des variations d'un objet (facilement)
    Problème
        -> Les autres font ça de telle manière dont voici les limites
    Solution
        exploiter :
        -> la formalisation des objets et des opérations
        -> la formalisation de détection d'événements topologiques

\section{speech}
\label{sec:org1e4df19}

Reevaluation mechanisms are, nowadays, made available in most professional grade
CAD softwares allowing users to generate variations of an object by editing 
operations' parameters. 

Nowadays, most existing methods require tracking numerous topological entities
and also consider the persistent naming problem only from the parameters'
 modifications of a parametric specification standpoint.

Unlike these methods, our approach is to extend the persistent naming problem 
to rule-based graph transformation modelling systems.

By using a formalism for operations, it is easier syntactically analyse rules
rather than symbolically analyse code.

This take allows us to propose:
\begin{itemize}
\item a persistent naming scheme representing unique topological parameter
histories
\item a reevaluation mechanism where only essential cells are being tracked
\item a support for both the edition of parameters and operations (that is,
the addition, deletion or displacement of operations) within a
parametric specification
\end{itemize}

In order to this, our method is able to represent the evolution of cells during
the initial evaluation of an object using Directed Acyclic Graphs. 

At reeavaluation, each directed acyclic graph is updated according to the changes
in order to match the topological parameters in the new object.

In addition, with such a mechanism reevaluation, we have the ability to 
setup custom reevaluation strategies such as, applying the same operation onto
each subdivisions of a cell.
\end{document}